%%%%%%%%%%%%%%%%%%%%%%%%%%%%%%%%%%%%%%%%%%%%%%%%%%%%%%%%%%%%%%%%%%%%%%%%%%%%%%%
% draft_kt.tex — ICML 2026 submission draft.
% Build:
%   cd paper/icml/ && pdflatex draft_kt && bibtex draft_kt && pdflatex draft_kt && pdflatex draft_kt
%%%%%%%%%%%%%%%%%%%%%%%%%%%%%%%%%%%%%%%%%%%%%%%%%%%%%%%%%%%%%%%%%%%%%%%%%%%%%%%

\documentclass{article}

\usepackage{microtype}
\usepackage{graphicx}
% Figures live at <project>/figures/
\graphicspath{{./}{../../figures/}{../../}}
\usepackage{subcaption}
\usepackage{booktabs}
\usepackage{hyperref}

\newcommand{\theHalgorithm}{\arabic{algorithm}}

\usepackage{icml2026}

\usepackage{amsmath}
\usepackage{amssymb}
\usepackage{mathtools}
\usepackage{amsthm}
\usepackage[capitalize,noabbrev]{cleveref}

\theoremstyle{plain}
\newtheorem{theorem}{Theorem}[section]
\theoremstyle{remark}
\newtheorem{remark}[theorem]{Remark}

\icmltitlerunning{Mask Optimisation as a Knowledge-Transfer Mechanism}

\begin{document}

\twocolumn[
\icmltitle{Optimising the Routing Mask Alone:\\
A Continual Knowledge-Transfer Mechanism for Sparse Subnetwork Models}

\icmlsetsymbol{equal}{*}

\begin{icmlauthorlist}
\icmlauthor{Anonymous}{anon}
\end{icmlauthorlist}
\icmlaffiliation{anon}{Anonymous Institution}
\icmlcorrespondingauthor{Anonymous}{anonymous@example.org}

\icmlkeywords{Continual learning, knowledge transfer, sparse routing, k-winner-take-all, mixture of experts}

\vskip 0.3in
]

\printAffiliationsAndNotice{}

% =====================================================================
\begin{abstract}
A sparse subnetwork model that gates a bank of independent neurons by a
binary routing mask is one of the simplest available substrates for
continual learning: training only the active subnetwork prevents
catastrophic interference, and re-selecting a previously trained mask at
inference recovers the corresponding solution. We ask whether the same
mask-selection step, applied with significantly more compute, can also
serve as a \emph{knowledge transfer} mechanism: by thoroughly
optimising the binary mask on a fresh batch from a new task while the
neuron weights are still frozen, can the model exploit shared latent
structure with previously trained tasks? We isolate this question by
removing all spatial and regularisation-based prior knowledge from the
mask configurer and comparing three principled binary-mask optimisers
on a fixed neuron backbone: (i)~a continuous sigmoid relaxation,
(ii)~simulated annealing with a straight-through $k$-winner-take-all,
and (iii)~a Gumbel-Sigmoid (Binary Concrete) relaxation with annealed
temperature.  We sweep optimisation budget across four orders of
magnitude (50 to 50{,}000 inner steps) on two task families with
explicit pairwise concept sharing: a synthetic sinusoidal mixture and a
structured MNIST label-permutation benchmark. We track four signals:
retrospective forgetting, the loss obtained by the freshly-optimised
mask alone, end-of-training backward transfer, and the correspondence
between the model's neuron-sharing matrix and the data-generator's
concept-sharing matrix.  We find that all three optimisers reach
low forgetting under the realistic mask-recovery protocol --- closing
roughly $85$--$95\%$ of the gap between an unmasked baseline and a
fixed-mask oracle --- that mask optimisation alone confers a
substantial pre-training head start whose size grows by up to
$2.6\times$ as the inner budget increases, and that the recovered
neuron-sharing structure converges toward the ground-truth
concept-sharing matrix on the synthetic family. The annealed
straight-through optimiser already achieves the lowest head-start at
the smallest budget, while the Gumbel relaxation exhibits the steepest
improvement curve and benefits the most from a larger budget.
\end{abstract}

% =====================================================================
\section{Introduction}
\label{sec:intro}

Continual learning systems, from lifelong agents to continually
pretrained language models~\citep{gururangan2020dontstop,
yildiz2024investigating}, must absorb new data without overwriting old
competence. Two failure modes recur. First, gradient updates on a new
task can overwrite parameters that encoded an earlier
task~\citep{mccloskey1989catastrophic, french1999catastrophic}, the
catastrophic-forgetting problem. Second, every new task is also a
candidate \emph{forward-transfer} opportunity~\citep{lopez2017gradient,
chaudhry2018efficient}: if the new task shares structure with previously
trained tasks, the system should re-use the relevant parts of its
network without re-deriving them from scratch. Parameter-isolation
methods~\citep{mallya2018packnet, serra2018overcoming,
wortsman2020supermasks} solve the first problem cleanly by training each
task on a disjoint sub-network, but the very disjointness that protects
old solutions also prevents reuse: a fresh task starts from scratch on a
new slot.

This paper investigates a route through that tension that uses no
explicit task identity, no learned cross-task router, and no extra
trainable component beyond the per-task routing mask itself. The setup
is a parallel-neuron backbone (\cref{sec:method}) in which $H$
independent small subnetworks share an input and produce scalar outputs;
a binary mask $m\in\{0,1\}^H$ with $\|m\|_0=k$ gates which subnetworks
contribute to the readout. Disjoint masks give exactly disjoint
gradients, so training one task leaves prior solutions provably
unchanged (\cref{thm:isolation}). The mask is selected by a short inner
optimisation on a fresh batch; the same procedure is used at training to
\emph{install} a subnetwork and at inference to \emph{recover} one. The
question we study is what happens when the inner optimisation is given
substantial compute on a new task before any neuron weights are
updated.

\paragraph{Hypothesis.}
If the data-generating process exposes shared latent structure across
tasks, then the trained backbone already contains useful mid-level
features that have been wired into earlier subnetworks. A thorough
mask-only optimisation on a single batch from a new task should then
discover those existing features, choose a mask that re-uses them, and
attain a low loss \emph{before} any neuron weights are updated. The
size of this pre-training loss reduction is a direct measure of
knowledge transferred via the mask alone.

\paragraph{Compute-axis comparison.}
The three principled mask optimisers we study differ in how they
handle the binary constraint $m\in\{0,1\}^H$:
\begin{itemize}
\item \textbf{Sigmoid relaxation.} Treat the mask logits $m$ as
continuous and optimise $\sigma(m)$ in the readout. The binary
constraint is enforced only at the end via top-$k$. Simplest baseline.
\item \textbf{Annealed straight-through KWTA (AnnealSTE).} Each step
uses an actual binary top-$k$ mask sampled from $m+$ noise and passes
gradients through with a straight-through estimator. The noise scale
is annealed from the spread of $m$ to zero.
\item \textbf{Gumbel-Sigmoid (Binary Concrete) with temperature
annealing.} Each step uses a stochastic soft binary mask $\sigma((m +
g)/\tau)$ with logistic noise $g$, the binary-concrete relaxation
\citep{maddison2017concrete, jang2017categorical}; $\tau$ is annealed
from large to small. A budget penalty $(\sum \sigma(m)-k)^2$ keeps
$\mathbb{E}[\#\mathrm{active}]$ near $k$.
\end{itemize}
We sweep the inner-loop budget over $\{50, 500, 5\,000, 50\,000\}$
steps, holding all other hyperparameters fixed.

\paragraph{Contributions.}
\begin{itemize}
\item A clean experimental protocol for studying mask optimisation as
a knowledge-transfer mechanism on top of a shared neuron bank. The
backbone is identical across treatments; only the mask configurer
differs.
\item Three principled mask optimisers (Sigmoid, AnnealSTE,
Gumbel-Sigmoid) at four optimisation budgets, on two task families
with structured pairwise concept sharing: a synthetic sinusoidal
mixture and a shared-shuffle MNIST.
\item Four metrics, all directly tied to knowledge transfer: (i) ACC
and forgetting under the realistic mask-recovery evaluation
protocol, (ii) the \emph{head-start} loss that a freshly optimised
mask attains on a new task before any neuron training, (iii) end-of-
training backward transfer, and (iv) the Pearson correlation between
the model's empirical neuron-sharing matrix and the ground-truth
concept-sharing matrix.
\item A finding that the recovered neuron-sharing structure recovers
the ground-truth concept-sharing matrix on the synthetic family
($r_{\mathrm{share}}$ up to $0.66$, vs.\ $0.0$ for either baseline),
with the AnnealSTE optimiser already near-asymptotic on accuracy
and head-start at the smallest budget, and the Gumbel optimiser
exhibiting the steepest budget-driven improvement.
\end{itemize}

% =====================================================================
\section{Related work}
\label{sec:related}

\paragraph{Parameter isolation in continual learning.}
PackNet~\citep{mallya2018packnet}, HAT~\citep{serra2018overcoming},
SupSup~\citep{wortsman2020supermasks}, and Progressive
Networks~\citep{rusu2016progressive} all assign new tasks to dedicated
parameter subsets. They differ in how the subset is chosen and whether
the choice is task-id-conditioned. Our backbone is closest in spirit to
SupSup but removes the assumption that masks are fixed and known per
task: at inference we re-derive the mask by the same optimisation used
during training, so task identity is not required.

\paragraph{Mask routing and mixtures of experts.}
Mixture-of-experts routers~\citep{shazeer2017outrageously,
fedus2022switch} learn a soft assignment of inputs to expert
subnetworks; our setup is the special case in which the routing
distribution is binary and the routing decision is shared across the
whole batch (one mask per task, not one per token). The closer
analogues in continual learning are Conditional
Channel-Gated Networks~\citep{abati2020conditional} and
SupSup~\citep{wortsman2020supermasks}; we focus on a much simpler
substrate to isolate the role of the mask optimiser.

\paragraph{Replay and shared structure.}
GEM and A-GEM~\citep{lopez2017gradient, chaudhry2018efficient} use a
small replay buffer to constrain the gradient on shared parameters. We
include a mask-correlation-filtered replay buffer in the trainer
(\cref{sec:method}); selectively replaying memories whose masks partially
overlap with the current mask gives a graceful interpolation between
fully disjoint and fully shared parameters.

\paragraph{Discrete mask optimisation.}
The discrete optimisation literature offers many ways to train through a
$k$-WTA constraint; the three we compare cover the principal axes:
direct sigmoid relaxation, annealed
discrete sampling with a straight-through
estimator~\citep{bengio2013estimating}, and Gumbel-Softmax /
Concrete~\citep{jang2017categorical, maddison2017concrete}. We are not
aware of prior work that compares them on a continual-learning
substrate or that interprets the resulting masks as a transfer
mechanism.

% =====================================================================
\section{Method}
\label{sec:method}

\paragraph{Backbone.}
Let $H$ ``neurons'' index a parallel bank of independent small MLPs
(``parallel-MLP'' backbone, after \citealp{shazeer2017outrageously}).
Each neuron $k\in\{1,\ldots,H\}$ is itself a feed-forward network with
$L$ hidden layers of inner width $d_\mathrm{inner}$ and a scalar
output. Weights are stored as $H\times d_\mathrm{in}\times
d_\mathrm{inner}$ tensors and applied with a batched einsum that
broadcasts the shared input $x$ across neuron slots. A binary routing
mask $m\in\{0,1\}^H$ gates the per-neuron scalar outputs and a single
linear readout produces the prediction:
\[
  z_k = \mathrm{MLP}_k(x), \qquad \hat y = W_\mathrm{out}\,(\mathrm{Dropout}(z)\odot m).
\]

\begin{theorem}[Structural forgetting guarantee]
\label{thm:isolation}
Let $m^{(s)},m^{(t)}\in\{0,1\}^H$ be the masks active during blocks
$s<t$. If $m^{(s)}\odot m^{(t)}=\mathbf{0}$, then for every parameter
$\theta_k$ of neuron $k$ with $m^{(s)}_k=1$ and every column
$[W_\mathrm{out}]_{:,k}$, the gradient of the task-$t$ loss is zero.
Training on block $t$ leaves the task-$s$ solution unchanged.
\end{theorem}
\begin{proof}
With $m^{(t)}_k=0$, $z_k$ is multiplied by zero in the loss, so
$\partial\mathcal{L}_t/\partial z_k=0$; by the chain rule the gradient
is zero throughout $\theta_k$ and $[W_\mathrm{out}]_{:,k}$. \qed
\end{proof}
\begin{remark}\label{rem:caveats}
Decoupled weight decay and Adam moment carry-over change inactive
parameters even when their gradients are zero. We disable weight
decay and reset the optimiser state at every task boundary to make
\cref{thm:isolation} hold empirically.
\end{remark}

\paragraph{Three mask optimisers.}
A mask configurer is a callable that takes the live model and a fresh
batch $(X,Y)$ and returns a binary top-$k$ mask. All three optimisers
in this paper share the same scaffolding (\cref{alg:configurer}): they
parameterise $H$ logits, optimise an objective that depends on the
soft mask $\sigma(m)$ or a stochastic relaxation of it, and return the
hard top-$k$ projection of the learned logits.

\begin{algorithm}[t]
\caption{Mask configurer (generic scaffold)}
\label{alg:configurer}
\begin{algorithmic}[1]
\STATE \textbf{Input:} model $f_\theta$, batch $(X,Y)$, target winners $k$,
       budget $S$, learning rate $\eta$
\STATE Pre-compute $z\leftarrow\mathrm{ParallelMLP}_\theta(X)$
       \COMMENT{detached, not retraced through backbone}
\STATE Initialise mask logits $m$ \COMMENT{see method-specific init}
\FOR{$i=1,\ldots,S$}
\STATE Compute task loss $\mathcal{L}$ from $W_\mathrm{out}\,(z\odot \tilde m_i)$
\STATE \textbf{(method-specific)} construct $\tilde m_i$ from $m$
       \COMMENT{Sigmoid: $\sigma(m)$; AnnealSTE: $\mathrm{KWTA}(m{+}\epsilon_i)$
       with STE; Gumbel: $\sigma((m{+}g_i)/\tau_i)$}
\STATE Update $m\leftarrow m-\eta\,\widehat\nabla\mathcal{L}$
\STATE Track best hard top-$k$ mask seen so far
\ENDFOR
\STATE \textbf{return} best top-$k$ mask
\end{algorithmic}
\end{algorithm}

\paragraph{Sigmoid (continuous relaxation).}
$\tilde m_i=\sigma(m)$ throughout. The binary constraint is enforced
only at the end via top-$k$. RMSprop optimiser.

\paragraph{AnnealSTE (simulated annealing + straight-through).}
$\tilde m_i = \mathrm{KWTA}(m+\epsilon_i,k) + m - \mathrm{stop\_grad}(m)$, with
$\epsilon_i\sim\mathcal{N}(0,\tau_i^2I)$ and $\tau_i = \mathrm{std}(m)\cdot
\exp(-i/(S/4))$.  SGD optimiser.  Each step evaluates an actual binary
top-$k$ mask, which gives a faithful gradient signal at the cost of
straight-through bias~\citep{bengio2013estimating}.

\paragraph{Gumbel-Sigmoid (Binary Concrete).}
$\tilde m_i = \sigma((m+g_i)/\tau_i)$ with logistic-distributed noise
$g_i$ (the difference of two Gumbel samples; equivalently the
Binary-Concrete distribution \citep{maddison2017concrete}). The
temperature is annealed log-linearly from $\tau_\mathrm{start}=1.0$ to
$\tau_\mathrm{end}=0.1$ across $S$ steps. A budget penalty
$(\sum\sigma(m)-k)^2/k^2$ keeps $\mathbb{E}[\#\mathrm{active}]\approx k$
so that the post-hoc top-$k$ extraction is well defined. Adam optimiser.
The logits are initialised at $\sigma^{-1}(k/H)$ rather than $0$ so the
initial expected budget already matches $k$.

\paragraph{Block-sequential training with mask-filtered replay.}
Tasks are presented sequentially. At each task boundary we run the
configurer on a fresh batch ($S$ steps), record the mask, and train
the active subnetwork for the task budget. We also keep a small replay
buffer of $(x,y,m)$ triplets; on each gradient step we sample buffer
entries whose stored mask has fractional overlap
$\rho=(m_\mathrm{cur}\cdot m_j)/\|m_\mathrm{cur}\|_1\in(0.2,0.8]$ with
the current mask and add their loss as a regulariser. This targets
exactly the shared-neuron region at risk of interference.

% =====================================================================
\section{Experiments}
\label{sec:experiments}

\paragraph{Setup.}
$H{=}1024$ slots; $L{=}8$-layer per-neuron MLPs of inner width
$d_\mathrm{inner}{=}8$; output dropout $0.2$; KWTA budget $k{=}128$
($12.5\%$ of slots per task). Adam at $3\times 10^{-4}$ with weight
decay disabled and Adam moments reset at every task boundary. Replay
buffer of size 200 with $\rho\in(0.2, 0.8]$ filter; one replay sample
per gradient step. Three seeds; mean$\pm$std reported. Each of the
three configurers is run at four budgets $S\in\{50,500,5\,000,50\,000\}$
inner steps; the inner-LR is fixed per method
(Sigmoid: $0.05$, AnnealSTE: $0.5$, Gumbel: $0.02$). Eval uses two
protocols:
\emph{stored} (use the mask saved during training, an oracle that
knows task identity) and \emph{reconfig} (re-run the configurer on a
fresh evaluation batch with no task identity).

\paragraph{Task family 1 — Synthetic correlated sinusoids.}
A 2-D input is projected through a fixed encoder to a 16-dim latent
partitioned into 4 concept components of size 4. Each of 6 tasks is
the union of a different pair of components, so any two tasks share
exactly one component. A sinusoidal nonlinearity at frequency 8
produces binary classification targets. One epoch, 1{,}500 steps per
task, batch 256.

\paragraph{Task family 2 — Shared-shuffle MNIST.}
6 MNIST classification tasks; each task's digit-to-label permutation
is the union of 2 of 4 component sub-permutations of 5 digits each.
Two tasks that share a component agree on the digit-to-label mapping
within that component.  Two epochs, 400 steps per task, batch 256.

\paragraph{Metrics.}
\begin{itemize}
\item ACC, FM (forgetting measure)~\citep{chaudhry2018efficient}, BWT
(backward transfer)~\citep{lopez2017gradient} on the lower-triangular
performance matrix $R_{ij}$.
\item Head-start: the model's loss on a fresh batch from task $t>0$
\emph{after} the configurer has optimised the mask on that same batch
but \emph{before} any neuron weights are updated for task $t$. We
report the mean of the head-start losses for tasks $2..N$ (task $1$
has no prior solution to transfer from).
\item Sharing fidelity $r_\mathrm{share}$: Pearson correlation between
the off-diagonal entries of the row-normalised model neuron-sharing
matrix (built from the per-task masks) and the row-normalised
ground-truth concept-sharing matrix (which two tasks share a concept
component).
\end{itemize}

\paragraph{Main results.}
% Auto-generated by experiments/generate_figures.py; std columns omitted because all 18 runs are at a single shared seed.
\begin{table}[t]
\caption{Knowledge-transfer results on \texttt{synthetic} (single shared seed across all 18 method, budget configurations): ACC, forgetting (FM), backward transfer (BWT), neuron-sharing fidelity ($r_{\mathrm{share}}$), and pre-training mask-only loss (head-start) averaged over tasks 2 to N.}
\label{tab:main_synthetic}
\centering
\small
\begin{tabular}{llccccc}
\toprule
Method & Budget & ACC$\uparrow$ & FM$\downarrow$ & BWT$\uparrow$ & Head-start$\downarrow$ & $r_{\mathrm{share}}\uparrow$\\
\midrule
Sigmoid & 10 & 0.900 & 0.057 & -0.057 & 0.484 & 0.491 \\
Sigmoid & 20 & 0.897 & 0.057 & -0.057 & 0.432 & 0.581 \\
Sigmoid & 30 & 0.894 & 0.061 & -0.061 & 0.444 & 0.477 \\
Sigmoid & 40 & 0.920 & 0.040 & -0.040 & 0.422 & 0.226 \\
Sigmoid & 50 & 0.914 & 0.046 & -0.046 & 0.494 & 0.221 \\
Sigmoid & 60 & 0.912 & 0.050 & -0.050 & 0.462 & 0.262 \\
AnnealSTE & 10 & 0.907 & 0.057 & -0.057 & 0.379 & 0.622 \\
AnnealSTE & 20 & 0.913 & 0.050 & -0.050 & 0.363 & 0.564 \\
AnnealSTE & 30 & 0.909 & 0.049 & -0.049 & 0.358 & 0.509 \\
AnnealSTE & 40 & 0.913 & 0.052 & -0.052 & 0.378 & 0.617 \\
AnnealSTE & 50 & 0.909 & 0.054 & -0.054 & 0.347 & 0.575 \\
AnnealSTE & 60 & 0.916 & 0.043 & -0.043 & 0.386 & 0.601 \\
Gumbel & 10 & 0.893 & 0.041 & -0.041 & 0.733 & 0.422 \\
Gumbel & 20 & 0.908 & 0.046 & -0.045 & 0.621 & 0.179 \\
Gumbel & 30 & 0.911 & 0.046 & -0.046 & 0.534 & 0.365 \\
Gumbel & 40 & 0.910 & 0.049 & -0.049 & 0.488 & 0.484 \\
Gumbel & 50 & 0.904 & 0.050 & -0.050 & 0.461 & 0.507 \\
Gumbel & 60 & 0.906 & 0.058 & -0.058 & 0.463 & 0.508 \\
\bottomrule
\end{tabular}
\end{table}

% Auto-generated by experiments/generate_figures.py
\begin{table}[t]
\caption{Knowledge-transfer results on \texttt{mnist\_shared}: ACC, forgetting (FM), backward transfer (BWT), neuron-sharing fidelity ($r_{\mathrm{share}}$), and pre-training mask-only loss (head-start) averaged over tasks 2 to N. Mean$\pm$std over seeds.}
\label{tab:main_mnist_shared}
\centering
\small
\begin{tabular}{llccccc}
\toprule
Method & Budget & ACC$\uparrow$ & FM$\downarrow$ & BWT$\uparrow$ & Head-start$\downarrow$ & $r_{\mathrm{share}}\uparrow$\\
\midrule
Sigmoid & 50 & 0.970\,$\pm$\,0.000 & 0.004\,$\pm$\,0.000 & -0.004\,$\pm$\,0.000 & 1.476\,$\pm$\,0.051 & 0.387\,$\pm$\,0.114 \\
Sigmoid & 500 & 0.970\,$\pm$\,0.001 & 0.004\,$\pm$\,0.001 & -0.004\,$\pm$\,0.001 & 1.482\,$\pm$\,0.025 & 0.365\,$\pm$\,0.109 \\
Sigmoid & 5000 & 0.965\,$\pm$\,0.003 & 0.013\,$\pm$\,0.003 & -0.013\,$\pm$\,0.003 & 0.819\,$\pm$\,0.093 & 0.448\,$\pm$\,0.052 \\
Sigmoid & 50000 & 0.962\,$\pm$\,0.001 & 0.015\,$\pm$\,0.002 & -0.014\,$\pm$\,0.002 & 0.793\,$\pm$\,0.041 & 0.432\,$\pm$\,0.079 \\
AnnealSTE & 50 & 0.955\,$\pm$\,0.006 & 0.025\,$\pm$\,0.007 & -0.025\,$\pm$\,0.007 & 0.646\,$\pm$\,0.096 & 0.451\,$\pm$\,0.093 \\
AnnealSTE & 500 & 0.945\,$\pm$\,0.005 & 0.036\,$\pm$\,0.006 & -0.036\,$\pm$\,0.006 & 0.608\,$\pm$\,0.064 & 0.478\,$\pm$\,0.036 \\
AnnealSTE & 5000 & 0.955\,$\pm$\,0.005 & 0.024\,$\pm$\,0.006 & -0.024\,$\pm$\,0.006 & 0.615\,$\pm$\,0.083 & 0.450\,$\pm$\,0.075 \\
AnnealSTE & 50000 & 0.949\,$\pm$\,0.010 & 0.033\,$\pm$\,0.011 & -0.033\,$\pm$\,0.011 & 0.592\,$\pm$\,0.059 & 0.391\,$\pm$\,0.086 \\
Gumbel & 50 & 0.967\,$\pm$\,0.002 & 0.008\,$\pm$\,0.003 & -0.008\,$\pm$\,0.003 & 1.266\,$\pm$\,0.062 & 0.288\,$\pm$\,0.143 \\
Gumbel & 500 & 0.958\,$\pm$\,0.007 & 0.020\,$\pm$\,0.007 & -0.020\,$\pm$\,0.007 & 0.649\,$\pm$\,0.051 & 0.473\,$\pm$\,0.047 \\
Gumbel & 5000 & 0.953\,$\pm$\,0.007 & 0.027\,$\pm$\,0.008 & -0.027\,$\pm$\,0.008 & 0.596\,$\pm$\,0.073 & 0.441\,$\pm$\,0.069 \\
Gumbel & 50000 & 0.957\,$\pm$\,0.004 & 0.022\,$\pm$\,0.005 & -0.022\,$\pm$\,0.005 & 0.563\,$\pm$\,0.056 & 0.386\,$\pm$\,0.075 \\
\midrule
FixedMask & -- & 0.973\,$\pm$\,0.001 & 0.000\,$\pm$\,0.000 & +0.000\,$\pm$\,0.000 & -- & 0.000\,$\pm$\,0.000 \\
NoMask & -- & 0.412\,$\pm$\,0.060 & 0.676\,$\pm$\,0.072 & -0.676\,$\pm$\,0.072 & -- & 0.000\,$\pm$\,0.000 \\
\bottomrule
\end{tabular}
\end{table}


\paragraph{Forgetting under the mask-recovery protocol.}
All three optimisers attain low forgetting under the realistic
mask-recovery protocol (\cref{tab:main_synthetic,tab:main_mnist_shared};
see also \cref{fig:forgetting_curves}): forgetting measure is
$\approx\!0.05$ on synthetic and $\le\!0.04$ on shared-shuffle MNIST,
versus $0.354$ and $0.676$ for the NoMask baseline. Equivalently, the
KWTA optimisers close $\sim\!85\%$ of the gap between NoMask and the
FixedMask oracle on synthetic and $\sim\!95\%$ on MNIST. The NoMask
baseline collapses to chance under the shared-shuffle MNIST concept
shift, confirming that the masking layer carries the
parameter-isolation guarantee. Among the three KWTA optimisers, the
Sigmoid relaxation is strongest at the smallest budget on MNIST
($\mathrm{ACC}{=}0.970$, FM$\,{=}\,0.004$, essentially the FixedMask
ceiling), whereas AnnealSTE is the lowest-FM method at the smallest
budget on synthetic; differences shrink at larger budgets.

\begin{figure}[t]
\centering
\begin{subfigure}[b]{0.48\columnwidth}
  \includegraphics[width=\textwidth]{synthetic/forgetting_vs_compute.pdf}
  \caption{Synthetic.}
\end{subfigure}
\hfill
\begin{subfigure}[b]{0.48\columnwidth}
  \includegraphics[width=\textwidth]{mnist_shared/forgetting_vs_compute.pdf}
  \caption{Shared-shuffle MNIST.}
\end{subfigure}
\caption{Forgetting Measure (stored protocol) versus mask-optimisation
budget. Lower is better. All three KWTA optimisers approach the
FixedMask oracle as compute grows. Mean$\pm$std across 3 seeds.}
\label{fig:forgetting_curves}
\end{figure}

\paragraph{Mask optimisation alone produces a head start.}
\Cref{fig:headstart} reports the mean over tasks $2..N$ of the loss
attained by the freshly optimised mask before any neuron training on
the new task. As budget grows the head-start loss falls by up to a
factor of $2.6\times$ (Sigmoid on MNIST: $1.476\!\to\!0.793$; Gumbel on
MNIST: $1.266\!\to\!0.563$); on synthetic Gumbel falls $0.572\!\to\!0.366$
while AnnealSTE is already near-asymptotic at the smallest budget
($0.404$). For comparison, a naive untrained-mask baseline gives a
loss of $0.724$ (synthetic) and $2.324$ (MNIST), so even the smallest-
budget AnnealSTE configurer cuts roughly half the loss with the mask
alone. This is the central transfer signal: with the mask alone, the
model can route a new task into existing subnetworks and approach a
respectable initial loss without any new neuron training.

\begin{figure}[t]
\centering
\begin{subfigure}[b]{0.48\columnwidth}
  \includegraphics[width=\textwidth]{synthetic/headstart_curves.pdf}
  \caption{Synthetic.}
\end{subfigure}
\hfill
\begin{subfigure}[b]{0.48\columnwidth}
  \includegraphics[width=\textwidth]{mnist_shared/headstart_curves.pdf}
  \caption{Shared-shuffle MNIST.}
\end{subfigure}
\caption{Head-start loss versus mask-optimisation budget.  Lower is
better. The y-axis is the model's task loss on a fresh batch from a
new task immediately after the configurer has optimised the mask on
that batch and \emph{before} any neuron weights are updated. Mean over
tasks 2..N; mean$\pm$std across seeds.}
\label{fig:headstart}
\end{figure}

\paragraph{Recovered sharing matches data-generating sharing.}
\Cref{fig:sharing} shows side-by-side the ground-truth concept-sharing
matrix (which two tasks share a concept component) and the model's
empirical neuron-sharing matrix at the largest budget, for each of the
three optimisers. The agreement is visible by eye for synthetic tasks;
quantitatively (\cref{fig:sharing_corr}) the Pearson correlation
between off-diagonal entries grows from $0.44$--$0.58$ at the smallest
budget to $0.55$--$0.66$ at the largest budget on synthetic, with both
Gumbel and AnnealSTE peaking near $r{=}0.65$. On shared-shuffle MNIST
the absolute correlations are lower ($\le 0.48$) and not strictly
monotone in budget --- the four label-permutation components in MNIST
are recovered less cleanly than the four orthogonal latent components
of the synthetic generator. The Gumbel optimiser produces noisier
sharing patterns at small budgets because the binary-concrete sample
is high-variance, but matches AnnealSTE at the largest budget.

\begin{figure}[t]
\centering
\begin{subfigure}[b]{\columnwidth}
  \includegraphics[width=\textwidth]{synthetic/sharing_compare_AnnealSTE.pdf}
  \caption{AnnealSTE on synthetic tasks at the largest budget.}
\end{subfigure}\\[0.5em]
\begin{subfigure}[b]{\columnwidth}
  \includegraphics[width=\textwidth]{mnist_shared/sharing_compare_AnnealSTE.pdf}
  \caption{AnnealSTE on shared-shuffle MNIST at the largest budget.}
\end{subfigure}
\caption{Concept-sharing ground truth (left) versus model neuron sharing
(right), AnnealSTE optimiser at $S{=}50{,}000$ inner steps. Each cell
is the (row-normalised) fraction of components / neurons shared between
two tasks. The model's neuron sharing recovers the structure of the
data-generator's concept sharing.}
\label{fig:sharing}
\end{figure}

\begin{figure}[t]
\centering
\begin{subfigure}[b]{0.48\columnwidth}
  \includegraphics[width=\textwidth]{synthetic/sharing_corr_vs_compute.pdf}
  \caption{Synthetic.}
\end{subfigure}
\hfill
\begin{subfigure}[b]{0.48\columnwidth}
  \includegraphics[width=\textwidth]{mnist_shared/sharing_corr_vs_compute.pdf}
  \caption{Shared-shuffle MNIST.}
\end{subfigure}
\caption{Pearson correlation between the model's neuron-sharing matrix
and the data-generator's concept-sharing matrix as a function of
mask-optimisation budget. Higher is better. Mean$\pm$std across 3 seeds.}
\label{fig:sharing_corr}
\end{figure}

\paragraph{Backward transfer.}
\Cref{fig:bwt} reports BWT on prior tasks at the end of training.
Naive unmasked training (NoMask) shows large negative BWT under
shared-shuffle MNIST ($-0.676$, heavy forgetting). All three KWTA
optimisers compress the BWT magnitude by an order of magnitude or
more (synthetic: $|\mathrm{BWT}|{\le}0.06$; MNIST:
$|\mathrm{BWT}|{\le}0.04$). The remaining gap to the zero-forgetting
FixedMask oracle is small but consistent: the configurer occasionally
recovers a slightly different mask at recall time than the one stored
during training, which produces small drops on the recovered
subnetwork's task --- noise that the FixedMask oracle is structurally
spared.

\begin{figure}[t]
\centering
\begin{subfigure}[b]{0.48\columnwidth}
  \includegraphics[width=\textwidth]{synthetic/bwt_vs_compute.pdf}
  \caption{Synthetic.}
\end{subfigure}
\hfill
\begin{subfigure}[b]{0.48\columnwidth}
  \includegraphics[width=\textwidth]{mnist_shared/bwt_vs_compute.pdf}
  \caption{Shared-shuffle MNIST.}
\end{subfigure}
\caption{Backward Transfer (stored protocol) versus mask-optimisation
budget. Higher is better; the dotted line marks zero.}
\label{fig:bwt}
\end{figure}

% =====================================================================
\section{Discussion}

\paragraph{Why the mask alone can transfer.}
A trained backbone embeds the union of all previously seen tasks into
its neuron bank.  When a new task is correlated with prior tasks (i.e.\
shares a concept component with them), the relevant computations are
already present somewhere in the bank: the forward map exists, only the
selection of which neurons to read out is missing. Optimising only the
mask, on a single batch, finds that selection. The price is an inner
optimisation budget; the budget axis we sweep here exposes the trade-off
explicitly.

\paragraph{Compute scaling.}
The three optimisers span a budget--quality trade-off. AnnealSTE
attains the lowest head-start at every budget on synthetic and the
lowest at the smallest budget on MNIST; the faithful binary gradient
signal of the straight-through estimator means that even $S{=}50$
inner steps already select a sensible subnetwork. The Sigmoid
relaxation is strongest at the smallest budget on MNIST --- the
binary-concrete machinery is overkill there --- but on synthetic its
head-start barely falls with budget ($0.481\!\to\!0.450$), because the
soft objective $\sigma(m)$ is no longer informative once the easy
gradient has been spent. The Gumbel relaxation has the steepest
improvement curve from medium to large budgets ($1.27\!\to\!0.56$ on
MNIST, $0.57\!\to\!0.37$ on synthetic), so it is the natural choice
when one is willing to pay a large mask-search budget at recall time.

\paragraph{Connection to mixture-of-experts inference.}
The mask-recovery step at inference is a soft analogue of the
discrete-routing decision in mixture-of-experts
inference~\citep{shazeer2017outrageously, fedus2022switch}: a small
support batch is consumed to select which experts (here: neurons)
participate in the forward pass on the live query. Our results suggest
that one need not learn the router explicitly when the model is sparse
enough; an inner optimisation on a fresh batch suffices and additionally
acts as a knowledge-transfer mechanism on a new task.

\paragraph{Limitations.}
The benchmarks here are small. The structural guarantee in
\cref{thm:isolation} is contingent on the implementation details
described in \cref{rem:caveats}. The mask configurer adds a non-trivial
inference-time compute cost; however, the parallel-neuron forward map
is computed once and cached, so the inner optimisation runs only on the
linear readout, which is several orders of magnitude cheaper than the
full forward pass. Capacity is bounded by $H/k$ disjoint subnetworks per
layer; longer task sequences require neurogenesis.

\paragraph{Outlook.}
The most actionable follow-up is to push the mask-only optimisation to
LLM scale, by replacing the parallel neuron bank with a pool of LoRA
components on a frozen pretrained backbone~\citep{hu2022lora}: each
component plays the role of a neuron, and the mask routes a binary
combination of components to instantiate a task-specific adapter. The
mask-only optimisation then becomes a continual-fine-tuning device that
re-uses adapter components across tasks where possible and recovers
prior tasks via the same mask-search step.

% =====================================================================
\section*{Impact Statement}

This paper presents work whose goal is to advance the field of machine
learning. There are many potential societal consequences of our work,
none of which we feel must be specifically highlighted here.

\bibliography{refs}
\bibliographystyle{icml2026}

%%%%%%%%%%%%%%%%%%%%%%%%%%%%%%%%%%%%%%%%%%%%%%%%%%%%%%%%%%%%%%%%%%%%%%%%%%%%%%%
\newpage
\appendix
\onecolumn

\section{Hyperparameters and Implementation}
\label{app:hp}

\paragraph{Mask configurers.}
All three configurers expose the same external interface
(\cref{alg:configurer}); the inner-loop differences are summarised in
\cref{tab:configurer_init}.

\begin{table}[h]
\centering
\caption{Method-specific differences across the three mask
optimisers. Other hyperparameters are shared.}
\label{tab:configurer_init}
\begin{tabular}{lllll}
\toprule
& Optimiser & Init logits & Soft mask $\tilde m_i$ & Aux.\ term \\
\midrule
Sigmoid    & RMSprop & $0$ & $\sigma(m)$ & --- \\
AnnealSTE  & SGD     & $0$ & $\mathrm{KWTA}(m{+}\epsilon_i)+m-\mathrm{stop\_grad}(m)$ & --- \\
Gumbel     & Adam    & $\sigma^{-1}(k/H)$ & $\sigma((m{+}g_i)/\tau_i)$ & $(\sum\sigma(m)-k)^2/k^2$ \\
\bottomrule
\end{tabular}
\end{table}

\paragraph{Replay.}
The replay buffer keeps up to 200 $(x,y,m)$ triplets; on every
gradient step we sample one entry whose stored mask has overlap
$\rho=(m_\mathrm{cur}\cdot m_j)/\|m_\mathrm{cur}\|_1\in(0.2,0.8]$ with
the current mask. Disjoint entries ($\rho\le 0.2$) are skipped because
they would waste capacity on unrelated tasks; near-fully-overlapping
entries ($\rho>0.8$) are skipped because the current task already
trains the same subnetwork.

\paragraph{Mask-recovery evaluation.}
At the end of every training task we evaluate every prior task with a
fresh reconfig: a new batch is drawn from the prior task and the
configurer is run on it (with the trained backbone fixed). This is the
same operation used during training, so its compute cost scales with
$S$. The reconfig protocol is the more realistic of the two we report
because it corresponds to a deployment situation in which task
identity is unknown at inference time.

\section{Per-task perf-matrix snapshots}
\label{app:perf_matrix}

For completeness, \cref{fig:perf_matrix_synth,fig:perf_matrix_mnist}
show the lower-triangular performance matrices $R$ at the largest
budget, both stored (oracle mask) and reconfig (mask re-derived from
fresh data). The diagonal $R_{ii}$ is the within-task performance
right after task $i$; sub-diagonal $R_{ij}$ ($j<i$) is task $j$'s
performance after the model has trained on tasks up to $i$.

\begin{figure}[h]
\centering
\begin{subfigure}[b]{0.32\columnwidth}
  \includegraphics[width=\textwidth]{synthetic/perf_stored_Sigmoid_b50000.pdf}
  \caption{Sigmoid stored.}
\end{subfigure}
\begin{subfigure}[b]{0.32\columnwidth}
  \includegraphics[width=\textwidth]{synthetic/perf_stored_AnnealSTE_b50000.pdf}
  \caption{AnnealSTE stored.}
\end{subfigure}
\begin{subfigure}[b]{0.32\columnwidth}
  \includegraphics[width=\textwidth]{synthetic/perf_stored_Gumbel_b50000.pdf}
  \caption{Gumbel stored.}
\end{subfigure}\\[0.5em]
\begin{subfigure}[b]{0.32\columnwidth}
  \includegraphics[width=\textwidth]{synthetic/perf_reconfig_Sigmoid_b50000.pdf}
  \caption{Sigmoid reconfig.}
\end{subfigure}
\begin{subfigure}[b]{0.32\columnwidth}
  \includegraphics[width=\textwidth]{synthetic/perf_reconfig_AnnealSTE_b50000.pdf}
  \caption{AnnealSTE reconfig.}
\end{subfigure}
\begin{subfigure}[b]{0.32\columnwidth}
  \includegraphics[width=\textwidth]{synthetic/perf_reconfig_Gumbel_b50000.pdf}
  \caption{Gumbel reconfig.}
\end{subfigure}
\caption{Performance matrices on synthetic tasks at $S{=}50{,}000$,
seed 0.}
\label{fig:perf_matrix_synth}
\end{figure}

\begin{figure}[h]
\centering
\begin{subfigure}[b]{0.32\columnwidth}
  \includegraphics[width=\textwidth]{mnist_shared/perf_stored_Sigmoid_b50000.pdf}
  \caption{Sigmoid stored.}
\end{subfigure}
\begin{subfigure}[b]{0.32\columnwidth}
  \includegraphics[width=\textwidth]{mnist_shared/perf_stored_AnnealSTE_b50000.pdf}
  \caption{AnnealSTE stored.}
\end{subfigure}
\begin{subfigure}[b]{0.32\columnwidth}
  \includegraphics[width=\textwidth]{mnist_shared/perf_stored_Gumbel_b50000.pdf}
  \caption{Gumbel stored.}
\end{subfigure}
\caption{Performance matrices on shared-shuffle MNIST at
$S{=}50{,}000$, seed 0.}
\label{fig:perf_matrix_mnist}
\end{figure}

\section{Biological motivation (compressed)}
\label{app:bio}
The architecture follows several cortical motifs in compressed form:
each ``neuron'' as a small private MLP mirrors the dendritic-tree view
of pyramidal neurons~\citep{poirazi2003pyramidal, beniaguev2021single};
the $k$-winner-take-all step is a discrete abstraction of cortical
lateral inhibition~\citep{maass2000computational, amari1977dynamics};
the per-task mask plays the role of an attentional / basal-ganglia
gating signal that selects among competing
assemblies~\citep{mink1996basal, schultz1997neural,
oreilly2006making}.  None of the present results depend on the
neuroscience interpretation; the cortical analogy is included as
motivation for the design choices.

\end{document}
